% status: FULL
% Chapter 13 of the second edition. Plan: docs/STRUCTURE.md. Rules: AUTHORING.md.
% Measurements: research/measurements/2026-09-26-m1-pass3.md, Chapter 13 (the lab's
% runs, Metal dispatch costs, llama.cpp's encoding of one decode step);
% research/measurements/2026-09-23-m1-part2.md (dispatch cost from PyTorch).
\chapter{How Kernels Are Written Now: CUDA, Triton, Tile DSLs and Graph Capture}\label{ch:kernel-languages}

\begin{ChapterIntent}
A decode step launches hundreds of small kernels, and at small batch the time
spent launching them can rival the time spent running them. This chapter
measures what a launch costs, piece by piece, on a real GPU; follows how
llama.cpp and vLLM stop paying for it --- few submissions, CUDA graphs captured
per batch size, fusion, compilers, megakernels; and surveys how production
kernels are written now: CUDA C++ with CUTLASS and CuTe, Triton, Python DSLs
such as CuTe DSL and TileLang, and NVIDIA's tile-level CUDA, with what each
abstraction hides and what it exposes.
\end{ChapterIntent}

\OpeningSection{Five hundred kernels per token}

When llama.cpp decodes a token of Llama~3.2~3B on the M1 of \Chref{roofline},
it runs a graph of \textbf{875 nodes} (measured; the log is in the record
\texttt{2026-09-26-m1-pass3}). Many of them cost nothing: views, reshapes and
permutations only relabel memory. Its Metal backend turns the rest into
\textbf{481 GPU operations} per token: 197 matrix--vector products, 57
normalizations, 56 rotations, 56 cache writes, 56 residual additions, 28
attention calls, 28 activations, and three small gathers and copies. Each
attention call is itself two kernels here, 32 partial passes over the cache
and a merge (\Chref{decode-attention}), so the GPU starts about 509 kernels for
every token --- at the 19.8 tokens per second of that run, ten thousand
kernels a second, each of which the host must describe and the GPU must start.

Starting one is not free. On an H100, Hazy Research measured about 2.1
microseconds to launch a kernel on a stream, and about 1.3 with CUDA graphs;
a Llama-1B forward pass in a conventional engine is around a hundred kernels
\cite{spector2025megakernel}. A hundred launches at two microseconds is
0.2~ms --- against a forward pass that streams 2.5~GB of BF16 weights in about
0.75~ms at the H100's full bandwidth (derived). At batch one on a small model,
starting the work can cost a quarter of doing it.

On the M1 the effect is easy to see from Python. A chain of 1,000 dependent
additions on 1,024 numbers, each addition its own PyTorch operation on the
GPU, took 6.5~ms; one addition doing the same arithmetic took 0.21~ms, so each
extra operation cost about 6~microseconds, none of it arithmetic. Compiled
with \texttt{torch.compile}, the chain became a single generated GPU kernel
containing all thousand additions, launched once. It took 0.21~ms --- the cost
of one operation --- after 2.6~s of compilation on first use (measured; record
\texttt{2026-09-23-m1-part2}). This chapter's lab takes Python out of the way
and asks the GPU directly. The same chain costs 160~microseconds per addition
when the host waits for each one, 31 when it only queues them, and 1.5 when all
thousand travel in one submission (measured).

The kernels themselves are the other half of the story. Every chapter of this
part derived an algorithm --- a tiled matrix product, an online softmax, a
split reduction, a dequantizing inner loop. Someone has to write those
algorithms for each generation of hardware, and in 2026 they are written in
more languages than ever. This chapter is about both halves: how engines stop
paying to launch kernels, and how the kernels are written.

\NapkinSection{Napkin math: launch overhead against step time}

If a step launches $n_k$ kernels at $t_{\text{launch}}$ each and its useful
work takes $t_{\text{work}}$, the overhead fraction is
\begin{equation}
f_{\text{launch}} = \frac{n_k\,t_{\text{launch}}}{n_k\,t_{\text{launch}} + t_{\text{work}}}.
\label{eq:launch-fraction}
\end{equation}
The work term shrinks with smaller models, faster GPUs and smaller batches;
the launch term does not. Four machines and models, put through the equation
(derived; the M1's kernel count measured from llama.cpp's log and its launch
cost from this chapter's lab, the H100's kernel counts at Hazy Research's rate
of about six kernels per layer, and Qwen3-8B's step times from
\Chref{napkin-math}):

\begin{center}\small
\begin{tabular}{@{}lrrrr@{}}
\toprule
Machine, model, batch & Kernels & Launching & Work & $f_{\text{launch}}$ \\
\midrule
M1, Llama~3.2~3B \texttt{Q4\_K\_M}, 1 & 509 & 0.66~ms & 50~ms & 1.3\% \\
H100, Llama-1B BF16, 1 & 100 & 0.21~ms & 0.75~ms & 22\% \\
H100, Qwen3-8B BF16, 1 & 220 & 0.46~ms & 5.9~ms & 7\% \\
H100, Qwen3-8B BF16, 32 & 220 & 0.46~ms & 12.9~ms & 3.5\% \\
\bottomrule
\end{tabular}
\end{center}

The M1 and the H100 sit at opposite corners. The M1's GPU is slow for its
launch cost: a 3B model's step takes fifty milliseconds, and five hundred
kernels at about 1.3~microseconds each are just over one percent of it. The
H100 is fast: a 1B model's weights stream through it in three quarters of a
millisecond, and a hundred launches are a fifth of the step. Batching dilutes
the launches, because a step at batch 32 launches the same kernels as a step
at batch one while taking more than twice as long. Launch overhead is a
batch-one, small-model, fast-GPU problem --- exactly the regime of interactive
use and of the draft models in speculative decoding
(\Chref{speculative-decoding}) --- and it grows with every generation whose
bandwidth grows faster than its launch latency falls. Hazy Research's
megakernel, which removes launches altogether, gained more over vLLM on a B200
than on an H100: 3.5 times against 2.5 (the authors' measurements)
\cite{spector2025megakernel}.

\section{What a launch costs}

A kernel launch is a message from the CPU to the GPU. In CUDA, the host calls
a launch function with a grid shape, a block shape and the kernel's arguments;
the driver packs them into a command and appends it to a \Term{stream}, a queue
that the GPU's front end consumes in order, and the call returns before the
kernel has run. The front end sets up the grid and hands thread blocks to
streaming multiprocessors as they free up; when every block has finished, the
next command in the stream may start. Metal, which llama.cpp uses on Apple
silicon, has the same pieces under other names: the host records
\emph{dispatches} through an encoder into a \emph{command buffer}, commits the
buffer to a \emph{command queue}, and may wait for it to complete.

\begin{center}\small
\begin{tabular}{@{}>{\raggedright\arraybackslash}p{0.31\linewidth}>{\raggedright\arraybackslash}p{0.33\linewidth}>{\raggedright\arraybackslash}p{0.3\linewidth}@{}}
\toprule
Idea & CUDA & Metal \\
\midrule
An ordered queue of GPU work & stream & command queue \\
One kernel launch & \verb|kernel<<<grid, block>>>(args)| & a dispatch through an encoder \\
Many launches, one submission & CUDA graph, recorded once & command buffer, recorded each time \\
Order between dependent kernels & stream order & serial encoder, or barriers \\
Wait for the GPU & \texttt{cudaStreamSynchronize} & \texttt{waitUntilCompleted} \\
\bottomrule
\end{tabular}
\end{center}

\noindent Metal also has indirect command buffers, which can be recorded once
and reused, closer to a CUDA graph; llama.cpp's Metal backend at
\texttt{d006858} does not use them.

Three different costs hide under the name \emph{launch overhead}, and they
have different remedies.

\begin{description}
\item[Issue cost.] The host's time to describe a launch: the framework's
bookkeeping (Python, PyTorch's dispatcher, shape checks) and the driver's
encoding of the command. It is paid on the CPU, and it hurts when the GPU
finishes kernels faster than the host can issue them. Recording the launches
once and replaying them, a CUDA graph, removes most of it.
\item[Gap cost.] The GPU's time between dependent kernels. The next kernel
cannot start until the last block of the previous one has finished and its
writes are visible; then its own blocks must be scheduled and issue their
first loads. Graphs do not remove it. Fusing the kernels does, and so does a
persistent kernel that schedules its own work.
\item[Synchronization cost.] A round trip: the host waits for the GPU to
finish, learns that it has, and only then sends more. The remedy is not to
wait --- to queue work ahead and synchronize only when the host needs a
result, which in decoding is once per token, to sample.
\end{description}

The lab separates the three on the M1's GPU. It applies a chain of 1,000
dependent additions, $x \leftarrow x + c_k$, to 1,024 floats in five ways ---
each addition in its own command buffer, waited for; each in its own command
buffer, queued without waiting; all as dispatches in one command buffer, in a
serial encoder or in a concurrent one with a barrier between them, as llama.cpp
does; and all in one fused kernel. It checks every result bit for bit against
the CPU, and it reads the GPU's own timestamps to separate time on the GPU from
time spent getting work to it. In microseconds per addition, two runs of five
repetitions each (measured, record \texttt{2026-09-26-m1-pass3}):

\begin{center}\small
\begin{tabular}{@{}lrrr@{}}
\toprule
How the chain is submitted & Host & GPU & Encoding \\
\midrule
A command buffer per addition, each waited for & 150--169 & 5.3 & --- \\
A command buffer per addition, one wait & 30.7--31.3 & 30.5--31.2 & --- \\
One command buffer, serial encoder & 1.49--1.50 & 1.17--1.18 & 0.14 \\
One command buffer, barriers (as llama.cpp) & 1.81--1.84 & 1.30--1.33 & 0.30 \\
One fused kernel & 0.23--0.25 & 0.074 & 0.004 \\
\bottomrule
\end{tabular}
\end{center}

\EngineFigure{ch13-submission}{Four ways to submit a chain of dependent
kernels (schematic). (a) Waiting after each kernel pays a round trip per
kernel. (b) Queuing without waiting removes the round trip but still pays for
every command buffer. (c) One submission for the whole chain --- one Metal
command buffer, or one CUDA graph replay --- leaves only the gap between
dependent kernels. (d) Fusion leaves one kernel. Right: costs measured on the
Apple M1 GPU by the chapter's lab.}{fig:ch13-submission}

Read the table from the top. Waiting after every command buffer costs about
160~microseconds per addition, and the GPU's timestamps say it was busy for 5
of them; the other 150 are the round trip, spent submitting the buffer,
scheduling it, noticing that it has finished and waking the host. Queuing the
buffers without waiting removes the round trip but not the buffer: each still
costs about 31~microseconds, and since the host's time and the GPU's span
agree, the buffers are processed one after another at that rate. A command
buffer is a heavy object on this machine, and an engine should submit few of
them. Put all thousand dispatches in one buffer and each costs 1.2
microseconds of GPU time --- the gap cost of a dependent kernel, since the
arithmetic on 1,024 floats takes nanoseconds --- plus 0.14 microseconds of host
time to encode. The concurrent encoder with barriers, which llama.cpp uses so
that independent operations may overlap, pays a little more for each dependent
one: 1.3~microseconds on the GPU and 0.3 to encode. Fuse the chain and the
whole thing is one dispatch. Its 74~microseconds of GPU time, spread over a
thousand additions, is less per addition than the gap cost of one of them, and
the 0.24~microseconds per addition that the host sees is mostly the chain's
single round trip.

The same ladder appears on a CPU, much lower. The lab's Rust part runs the
chain on a worker thread that owns the buffer and takes commands from a queue,
as a GPU does from a stream. At 1,024 floats a round trip --- send an addition,
wait for the worker to report it done --- costs 2.7~microseconds per addition;
queuing each addition costs 0.073, a recorded chain sent as one command 0.058,
and a fused kernel 0.026 (measured). A thread in the same process is a much
nearer device than a GPU behind a driver: its round trip is sixty times cheaper
and a queued command four hundred times cheaper. But the order is the same,
and so is the lesson: waiting is dearest, submitting is next, and only fusion
removes work rather than overhead.

\EngineFigure{ch13-costs}{Cost per addition of a chain of 1,000 dependent
additions on 1,024 floats, by how the work is submitted, on a logarithmic
scale (measured). The GPU's ladder (Metal, host wall clock, mean of two runs)
and the CPU's (a worker thread fed by a queue) have the same order, three rungs
apart in places: a GPU is a remote device.}{fig:ch13-costs}

Larger buffers show where the fixed costs go. At a million floats a dispatch in
one command buffer costs 50 to 72~microseconds of GPU time, because each
addition now reads and writes 8~MB. The overheads shrink against the work:
waiting after each addition costs five to seven times as much as a dispatch in
one buffer, down from a hundred times, and a command buffer per addition 1.4
times as much; on the CPU the three unfused methods agree within a tenth. Only
fusion still changes the picture, applying the thousand additions at
6.2~microseconds each on the GPU, because it reads and writes the buffer once
instead of a thousand times (measured). Launch overhead and memory traffic are
different problems. A CUDA graph removes the first; only fusion removes the
second.

\section{How llama.cpp submits a token}

llama.cpp's answers to the three costs are visible in its source at commit
\texttt{d006858} and in the log of the decode step that opened this chapter.

\emph{Build the graph once.} Before each step the context checks whether the
new batch's graph would have the same topology as the last one. For decoding
it does, and the old graph is reused without being rebuilt or re-allocated
(\texttt{llama\_context::process\_ubatch}; the run reports ``graphs reused'',
and the environment variable \texttt{LLAMA\_GRAPH\_REUSE\_DISABLE} turns it
off). This is not a recording of GPU work, only of the host's description of
it, but it removes the cost of constructing 875 nodes for every token.

\emph{Few command buffers, encoded in parallel.} The Metal backend encodes a
token's graph into two command buffers. The calling thread encodes the first
$\max(64, n/10)$ of the $n$ nodes and commits them at once, so that the GPU
starts working, while a second thread encodes the rest into the other buffer;
a comment in the source records that one or two extra threads worked best on
an M1 Pro and an M2 Ultra (\texttt{ggml-metal-context.m}). At about
31~microseconds per command buffer on the M1, two per token cost 0.06~ms of a
50~ms step.

\emph{Barriers only where needed.} The encoder is of Metal's concurrent type,
in which dispatches may overlap. Before each operation the backend compares the
memory ranges it reads and writes with those of the operations since the last
barrier. If it writes nothing they touch and reads nothing they write, it runs
concurrently; otherwise the backend inserts a memory barrier and starts a new
set (\texttt{ggml-metal-ops.cpp}). In the decode step, 116 of the 481
operations ran without a barrier --- the key and value projections alongside
the query projection, the up projection alongside the gate --- and 365 waited
for the ones before them.

\emph{Fuse the obvious neighbours.} A normalization followed by a
multiplication by its weight, and an addition if one follows, becomes one
kernel; so does a chain of up to eight additions of the same shape. In the
decode step all 57 normalizations were fused with their weights, removing 57
dispatches and 57 trips of an activation vector through memory. The
environment variables \texttt{GGML\_METAL\_FUSION\_DISABLE} and
\texttt{GGML\_METAL\_CONCURRENCY\_DISABLE} turn the last two ideas off, which is
the quickest way to measure what they are worth on a given machine.

\emph{On NVIDIA GPUs, graphs.} The CUDA backend uses CUDA graphs on Ampere and
newer GPUs. A graph's first call runs eagerly; a second call whose nodes have
the same properties completes a warm-up and is captured; later calls with
unchanged properties replay the capture. When the properties change, the
warm-up starts again, and the executable graph is updated in place with
\texttt{cudaGraphExecUpdate} or, if that fails, instantiated anew. Graphs are
disabled for weights split across GPUs and for mixture-of-experts products
that would have to synchronize the stream (\texttt{ggml-cuda.cu}).

\section{Launch overhead and CUDA graphs}

A CUDA graph records a sequence of kernel launches once and replays it with a
single submission, removing most of the issue cost of every launch after the
first; Hazy Research's 2.1~microseconds fell to 1.3, and the gap cost remains.
The catch is that a graph captures fixed shapes and fixed memory addresses:
the replay launches the same kernels with the same arguments on the same
buffers. An engine that wants graphs therefore keeps the tensors a graph
touches at fixed addresses, copies each step's inputs into them, and captures
one graph per batch size it expects.

vLLM's rule is explicit in its source at commit \texttt{bcdacfc}
(\texttt{vllm/config/vllm.py}, \texttt{vllm/v1/cudagraph\_dispatcher.py}). It
captures batches of 1, 2 and 4 tokens, every multiple of 8 below 256, and every
multiple of 16 from 256 up to a maximum: 512 on most GPUs and 1,024 on
data-centre Blackwell, but never more than twice the number of sequences the
scheduler may run at once. A batch runs at the smallest captured size that
holds it, padded with dummy tokens, and a batch larger than the maximum runs
without a graph. With the H100's default limit of 1,024 sequences that is 51
sizes; on a B200, 83. The default mode, \texttt{FULL\_AND\_PIECEWISE}, captures
two kinds of graph for each size: a full graph, attention included, for batches
in which every request decodes one token, and piecewise graphs --- everything
except attention --- for batches that mix prefills and decodes, whose attention
work depends on each request's lengths \cite{vllm2026cudagraphs}. The graphs
are captured at startup, largest first, before the server accepts a request.

Padding is the price of a finite set of sizes, and the lab's third part
computes it. If every batch size from 1 to 512 were equally likely, vLLM's 51
sizes would waste 3.3\% of the computed slots on average; the worst case is a
batch of 9 padded to 16, 44\% padding. Ten graphs at the powers of two would
waste 24.5\% on average and up to 49.8\%, for a batch of 257 run at 512
(derived). The wasted slots cost less than their share suggests in decode,
where a step is limited by reading weights and cache and a padded token adds
arithmetic the GPU had idle anyway (\Chref{napkin-math}); they cost their full
share when a step is limited by compute, as a step carrying prefill chunks is
(\Chref{prefill-decode}). That is why the sizes are coarse only at the bottom
--- 1, 2, 4 and 8 --- and step by 8 and 16 above it, where token counts are
large: there padding costs its full share, the steps are small, and each extra
size is worth its memory and its capture time.

\section{Fusion and compilers}

Graphs remove the cost of launching kernels; fusion removes kernels, and with
them the memory traffic between them. The lab's second part measures the
traffic on the CPU with the three nodes llama.cpp fuses: RMSNorm, a
multiplication by the norm's weight and an addition, over 4,096 tokens of
width 3,072, 50~MB of activations. As three kernels the activations cross
memory three times each way; fused, once each way, so a pass limited by memory
bandwidth should be three times faster. It was 2.6 to 3.3 times faster in four
measurements (one and four threads, two runs), with the unfused passes moving
38 to 55~GB/s by the byte count, most of the M1's 67--68~GB/s
(\Chref{roofline}), and the outputs were identical to the bit, because the fused kernel rounds each
element exactly as the three passes do (measured). Here is the fused kernel;
the unfused version is the same loop split in three:

\CodeLines{../code/labs/ch13-launch/src/main.rs}{306}{317}

Fusion is not free. A fused kernel holds more state per thread and does more
work between memory accesses; when its arithmetic has too little parallelism
to hide its own latency, it can be slower than the kernels it replaced. The
lab shows this by shrinking the fused chain's tile from 64 elements to 8: its
additions then form too few independent chains to keep the CPU's adders busy,
and the fused chain became slower than a thousand separate passes, 0.61
microseconds per addition against 0.31 at 4,096 floats and 118 against 76 at a
million (measured). On a GPU the same trade appears as registers: a fused
kernel that needs more registers per thread runs fewer threads at once, and
may no longer hide memory latency (\Chref{matmul-hardware}).

Writing fused kernels by hand for every pair of neighbours does not scale, so
engines hand the job to compilers. vLLM compiles each model's forward pass
with \texttt{torch.compile}, wraps attention as an opaque operation, splits
the graph at attention, and compiles the pieces between --- the pieces its
piecewise CUDA graphs capture --- all before it serves a request, so no request
waits for a compiler \cite{vllm2025torchcompile}. The M1 measurement in the
opening is the same mechanism at small scale: PyTorch's compiler turned a
thousand operations into one generated kernel, at the price of 2.6~s of
compilation on first use.

The limiting case is a single \Term{megakernel}: the whole forward pass as one
persistent kernel that schedules its own work. Hazy Research built one for
Llama-1B and named three wastes it removes (the authors' measurements, May
2025). Launches, which even graphs leave at 1.3~microseconds each. Tails, the
idle multiprocessors at the end of a kernel whose blocks do not divide evenly
among them: 512 blocks on a B200's 148 multiprocessors leave 80 idle in the
last wave. And stalls, the thousands of cycles each kernel spends loading its
first weights, which could have overlapped the previous kernel's work. Their
megakernel runs an interpreter on every multiprocessor, executing a
precomputed list of instructions; it divides shared memory into thirteen
16~KiB pages that one instruction can hand to the next, so the next weights
start loading while the current computation finishes; and it replaces kernel
boundaries with counters in global memory that instructions increment and
wait on. It used 78\% of an H100's memory bandwidth where engines of the time
used at most half, and finished a forward pass in under a millisecond on an
H100 and under 680~microseconds on a B200 \cite{spector2025megakernel}. MPK
compiles models into such kernels automatically, from a graph of tasks at the
granularity of single multiprocessors executed by a runtime inside the kernel,
and reports up to 1.7 times lower end-to-end latency than serving systems that
launch a kernel per operator \cite{cheng2025mpk}.

\section{The kernel author's choices}

A kernel is a function run by thousands of threads at once. Writing one means
choosing, for a given operation and machine, how to divide the output among
thread blocks, how to stage inputs through shared memory and registers, which
matrix instructions to use, what precision to accumulate in, and how to
overlap loading the next tile with computing the current one
(\Chref{matmul-hardware}). Kernel languages differ in which of these choices
they make the programmer spell out and which they make for them. The more a
language decides, the faster kernels are written; the less it decides, the
closer they can get to the hardware's limits.

\section{One kernel, four ways}

The lab's kernel is the smallest possible: add a constant to every element of
an array. Written for four programming models, it shows where each puts the
work of mapping data to threads. (The CUDA, Triton and cuTile versions below
are sketches written from each language's documentation; this book's tests
cannot run them without an NVIDIA GPU. The Metal version is the lab's.) In
CUDA C++ the unit is one thread:

\begin{verbatim}
__global__ void add_one(float *x, float c, int n) {
    int i = blockIdx.x * blockDim.x + threadIdx.x;   // this thread's element
    if (i < n) x[i] += c;                            // the last block overhangs
}
// launch: add_one<<<(n + 255) / 256, 256>>>(x, c, n);
\end{verbatim}

\noindent The programmer computes each thread's index from its block and thread
numbers, chooses 256 threads per block, and guards the elements past the end.
Metal on Apple GPUs is the same model. The lab launches exactly as many threads
as there are elements, so its kernel needs no guard:

\CodeLines{../code/labs/ch13-launch/metal/kernels.metal}{8}{12}

\noindent In Triton the unit is a \emph{program} that owns a block of elements:

\begin{verbatim}
@triton.jit
def add_one(x_ptr, c, n, BLOCK: tl.constexpr):
    offs = tl.program_id(0) * BLOCK + tl.arange(0, BLOCK)   # this program's block
    mask = offs < n
    x = tl.load(x_ptr + offs, mask=mask)
    tl.store(x_ptr + offs, x + c, mask=mask)
# launch: add_one[(triton.cdiv(n, 256),)](x, c, n, BLOCK=256)
\end{verbatim}

\noindent Loads and stores move whole blocks, with a mask for the overhang; how
the block is spread over the threads of a warp, and in what widths memory is
read, is the compiler's choice. In cuTile Python the unit is a \emph{tile} of
an array, and the pointer arithmetic disappears:

\begin{verbatim}
@ct.kernel
def add_one(x, c, tile: ct.Constant[int]):
    i = ct.bid(0)                                           # this block's tile
    t = ct.load(x, index=(i,), shape=(tile,))
    ct.store(x, index=(i,), tile=t + c)
# launch: ct.launch(stream, ((n + 255) // 256, 1, 1), add_one, (x, c, 256))
\end{verbatim}

\noindent The kernel receives arrays that know their shapes, an index names a
tile rather than an element, and the edges are the runtime's problem: a load
that runs past the end of the array pads the missing elements, with a value
the kernel may choose, and a store ignores them. Tile shapes must be known when the kernel is compiled, and each
dimension must be a power of two \cite{nvidia2026cutile}.

For one addition the four versions are the same program. They part ways as
soon as an operation needs threads to cooperate. An RMSNorm row needs a sum
across its 3,072 elements. In CUDA or Metal the programmer writes it: each
thread sums a slice, the threads of a warp combine their partial sums with
shuffle instructions, and the warps combine theirs through shared memory, with
the programmer choosing the number of threads and how the slices interleave.
In Triton it is \texttt{tl.sum(x * x, axis=0)} over the block, and the compiler
writes the shuffles; in cuTile it is a reduction over a tile. A matrix product
goes further. In Triton it is \texttt{tl.dot} on two blocks, and the compiler
picks the tensor-core instructions and stages the operands through shared
memory, with \texttt{num\_stages} (3 by default) setting how many loads the
pipeline keeps in flight and \texttt{num\_warps} (4 by default) how many warps
share the work \cite{triton2026config}; in CUDA C++ every one of those choices
is the programmer's.

\begin{center}\small
\begin{tabular}{@{}>{\raggedright\arraybackslash}p{0.2\linewidth}>{\raggedright\arraybackslash}p{0.17\linewidth}>{\raggedright\arraybackslash}p{0.19\linewidth}>{\raggedright\arraybackslash}p{0.16\linewidth}>{\raggedright\arraybackslash}p{0.16\linewidth}@{}}
\toprule
Who decides & CUDA C++, Metal & CUTLASS and CuTe & Triton & cuTile \\
\midrule
Unit of code & a thread & a thread, with tensors described by layouts & a program owning a block & a block owning tiles \\
Mapping elements to threads & programmer & programmer, through layouts & compiler & compiler \\
Staging through shared memory & programmer & library templates & compiler & compiler \\
Matrix instructions & programmer & templates per architecture & compiler & compiler \\
Edges of arrays & explicit guards & predication & masks & padded loads, ignored stores \\
\bottomrule
\end{tabular}
\end{center}

\section{CUDA C++, CUTLASS and CuTe}

At the bottom is CUDA C++ with inline PTX: every choice explicit, every new
instruction available on the day the hardware ships. NVIDIA's CUTLASS library
packages the recurring structure of matrix kernels as C++ templates, and its
CuTe layer describes tensors and thread arrangements as composable layouts ---
the vocabulary in which Hopper's warpgroup MMA and TMA copies and Blackwell's
tensor-memory instructions are exposed (\Chref{matmul-hardware}). Most of the
fastest GEMM and attention kernels in production engines are CUTLASS or
CuTe underneath. The price is template-heavy code that is slow to compile and
hard to read.

CuTe's central object is the \Term{layout}: a shape and a stride, printed as
\texttt{Shape:Stride}, which map a coordinate to an offset by the inner product
of the coordinate with the stride \cite{nvidia2026cutelayout}. A $4\times2$
tile stored by rows is \texttt{(4,2):(2,1)}, and element $(i, j)$ lives at
offset $2i + j$; stored by columns it is \texttt{(4,2):(1,4)}, offset $i + 4j$.
Because shapes and strides can nest, one layout can say how a tile is divided
among the 128 threads of a Hopper warpgroup in the pattern its matrix
instruction requires, or how shared memory is permuted to avoid bank
conflicts; and because layouts compose, a kernel is written as a chain of
layout transformations from the tensor in global memory to each thread's
registers. The same algebra serves every architecture; what changes between
generations is which layouts the hardware's instructions accept.

\section{Triton: tiles in Python}

Triton, introduced in 2019, raised the unit of programming from the thread to
the tile: the programmer writes what one block computes on blocks of data, and
the compiler decides how threads cooperate within it \cite{tillet2019triton}.
It became the default way to write custom kernels in the PyTorch ecosystem,
and PyTorch's own compiler generates Triton code. Triton kernels are usually
autotuned: the author lists candidate block sizes, warp counts and pipeline
depths, and Triton times each on the first call for a new problem shape and
keeps the fastest. vLLM ships a Triton attention backend that runs where vendor
kernels are not available \cite{vllm2026attentionbackends}.

\section{CuTe DSL, TileLang and ThunderKittens}

A newer generation keeps Triton's productivity but hands explicit control back.
CuTe DSL exposes CuTe's layouts and hardware instructions from Python, with
just-in-time compilation. FlashAttention-4 is written in it, and its authors
report compile times 20--30 times shorter than the equivalent C++ templates
\cite{zadouri2026flashattention4}. TileLang separates what a kernel computes
from how it is scheduled \cite{wang2025tilelang}; ThunderKittens builds kernels
from $16\times16$ tiles and matches cuBLAS and FlashAttention-3 on H100 in its
authors' measurements \cite{spector2024thunderkittens}.

\section{CUDA Tile and cuTile}

In December 2025 NVIDIA added a tile-based programming model to CUDA itself:
CUDA Tile, with a virtual instruction set for tile programs (CUDA Tile IR) and
cuTile Python, a DSL on top of it \cite{nvidia2025cuda131}. Its reach has grown
release by release. The CUDA Tile release notes date support for Ampere and
Ada to CUDA 13.2 in March 2026 \cite{bentz2026cuda132}, Hopper to the release
aligned with CUDA 13.3 in May, and a developer preview for Rubin to September
\cite{nvidia2026cudatile}. Its significance for engines is portability: a tile
program written once can target tensor cores across generations, where a
CUTLASS kernel is typically rewritten per architecture.

\section{Beyond NVIDIA}

The same layering exists elsewhere, with different names: Metal shaders on
Apple silicon, used by llama.cpp; HIP and ROCm on AMD, where vLLM has its own
attention backends; and portable CPU kernels, which ggml provides for every
major instruction set (Appendix~\ref{app:cpu}). An engine that serves more
than one vendor's hardware carries a kernel set per backend, and the sets are
large: at commit \texttt{d006858} llama.cpp's Metal shader library defines 109
kernel functions, instantiated as 560 named kernels for different types and
shapes, and its CUDA backend has 61 kernel source files plus 116 generated
instantiations of templated ones. Every one needs correctness tests that
compare it with a reference (\Chref{fast-wrong-answers}).

\LedgerSection

Capturing decode steps as CUDA graphs (or submitting them in few command
buffers) and fusing the operations between attention calls:

\begin{ConstraintLedger}
Latency & buys & Launch gaps and round trips disappear: about a fifth of a 1B model's batch-one step on an H100, about a percent of a 3B model's on an M1 (derived). \\
Memory bandwidth & buys & Fused operations read and write their activations once, not once per operation: three times faster for norm, scale and shift in the lab (measured). \\
Compute & buys & The GPU stops idling between kernels. \\
Memory capacity & spends & One graph and its fixed buffers per captured size; vLLM captures 51 sizes on an H100 by default. \\
Latency (startup) & spends & Compilation and capture before the first request (2.6~s to compile a 1,000-operation chain on the M1); cacheable. \\
Correctness & risks & Padding to captured sizes and replayed addresses are new places for bugs; a fused kernel must round as the unfused ones did, or its outputs change. \\
\end{ConstraintLedger}

\LabSection{measure what a launch costs}

The lab in \texttt{code/labs/ch13-launch/} runs on any laptop CPU. A worker
thread plays the GPU: it owns a buffer and executes commands from a queue in
order, and the host can only send commands and wait on a fence:

\CodeLines{../code/labs/ch13-launch/src/main.rs}{116}{131}

\noindent The four ways of applying the chain differ only in what the host sends
and when it waits:

\CodeLines{../code/labs/ch13-launch/src/main.rs}{188}{211}

\noindent and the fused kernel keeps a tile of 64 elements in registers while
it applies every addition, in the same order as the separate passes:

\CodeLines{../code/labs/ch13-launch/src/main.rs}{72}{87}

\noindent Every method must reproduce, bit for bit, an oracle that adds the
constants to each element in turn on the host; only then is it timed. The
second part times RMSNorm, scale and shift as three passes and as one over
50~MB of activations, checks both against float64 and against each other, and
counts bytes; the third derives the padding that captured sizes cost. Run it
from \texttt{code/labs} with \texttt{cargo run --release -p ch13-launch}.

On a Mac, \texttt{metal/launch.m} is the GPU extension behind this chapter's
table of Metal costs. It runs the same chain on the GPU five ways, with the two
kernels of \texttt{metal/kernels.metal} compiled in Metal's safe-math mode so
that the GPU rounds every addition as the CPU does, and it reads the GPU's
timestamps for every command buffer; the README shows how to build it with
\texttt{xcrun clang}. Before running either program, predict: on your machine,
is a round trip closer to a microsecond or a millisecond, and how many
dispatches in one command buffer cost as much as one round trip? Then use your
own numbers in equation~\eqref{eq:launch-fraction} for the smallest model you
serve.

\HappenedSection{measuring the Metal round trip instead of guessing it}

Hermon's native attention kernel for Apple GPUs was correct --- it matched the
CPU path to $4.8\times10^{-8}$ --- and slower than the CPU below about 4,000
tokens of context on an M1. Every call cost about 0.94~ms however little work
it did. Hermon's developers first wrote down three fixes, based on a guessed
breakdown of that cost, and then measured the breakdown
(Hermon \texttt{docs/KERNEL\_DESIGN.md}, section K4.1):
\begin{itemize}
\item submitting a trivial kernel and waiting for it: \textbf{0.190~ms},
irreducible;
\item allocating three staging buffers: 0.009~ms, about 1\% --- not the lever
the guess assumed;
\item each extra dispatch sharing one command buffer: 0.008~ms.
\end{itemize}
Batching all 28 layers' attention into one command buffer would have removed
most of the overhead, and it was impossible: the projections and FFN between
attention calls run on the CPU through llama.cpp, so each layer's attention
result must come back before the next layer can start. The fix that shipped
was a gate: the Metal kernel runs only for shapes large enough to repay its
fixed cost, and the gate depends on the shape alone so that a request always
takes the same path: about 16 million elements of query tokens $\times$
heads $\times$ context $\times$ head width, calibrated on the M1 and declared
not to transfer to other devices (\texttt{csrc/attn.c} at commit
\texttt{2a3fd52}). The whole native kernel path is opt-in through an
environment variable, so PREVIEW. Two of three planned optimizations would have
optimized the wrong thing.

This chapter's lab measured the same round trip independently, with different
code on the same kind of machine: 150 to 169~microseconds for a command buffer
committed and waited for, of which the GPU was busy for about 5. Hermon's
0.19~ms was the platform's price, not a flaw in Hermon's code --- and the
lesson generalizes. A kernel that must return its result to the host before the
next one can start pays the round trip every time, whatever its own speed;
only moving the work on either side of it to the same device, as llama.cpp
does by running whole layers on the GPU, removes it.

\FrontierSection{2026-09-27}

Kernel authoring moved up a level in 2025--26. FlashAttention-4 is written in
CuTe DSL \cite{zadouri2026flashattention4}; in vLLM at commit \texttt{bcdacfc}
it is the version that the FlashAttention backend selects on Blackwell
(compute capability 10) when supported, while FlashInfer is the first choice
for causal attention there (\texttt{vllm/platforms/cuda.py} and
\texttt{fa\_utils.py}). NVIDIA made tiles a first-class CUDA programming model
and widened it every few months, to Ampere, Ada, Hopper and, in preview, Rubin
\cite{nvidia2026cudatile}. Megakernels moved from a single demonstration to a
research line: MPK compiles models into them automatically
\cite{cheng2025mpk}, and in September 2026 ForgeMegakernel used coding agents,
checked by a float64 oracle, to generate a decode megakernel per model,
reporting geometric-mean speedups of 1.21$\times$ over SGLang and 1.54$\times$
over MPK on one H100 \cite{li2026forgemegakernel} (the authors'
measurements). The constant through all of it is the constraint of this
chapter's napkin math: at small batch, the time between kernels is as real as
the time in them.

\ExercisesSection

\begin{enumerate}
\item \emph{Derivation.} Using equation~\eqref{eq:launch-fraction}, find the
batch size at which launch overhead falls below 2\% for an 8B model with 220
kernels per step at 2~$\mu$s each on an H100, taking the step time from
\Chref{napkin-math}.
\item \emph{Prediction.} A CUDA graph is captured for batch sizes 1, 2, 4, 8,
16 and 32. Predict the throughput cost of a steady load of 17 concurrent
requests, and propose a better set of sizes; compare yours with vLLM's rule.
\item \emph{Implementation.} Write a fused RMSNorm-plus-matrix--vector kernel
in Triton or a CPU equivalent in Rust, check it against separate operations,
and measure what fusion saves at batch one.
\item \emph{Falsification.} Find an operation that becomes slower when fused
with its neighbour, and explain which resource the fusion spent.
\item \emph{Measurement.} Run the lab's Metal extension (or, on an NVIDIA GPU,
time the same chain with CUDA events, eagerly and as a replayed CUDA graph) and
fill in this chapter's table of costs for your machine. Which of the three costs
dominates, and which remedy would you buy first?
\end{enumerate}

\section{Worked problems: padded batches, what graphs buy, and when fusion pays}

\subsection{The price of padding}
\textbf{Problem.} vLLM captures CUDA graphs on an H100 at batch sizes of 1, 2
and 4 tokens, every multiple of 8 below 256 and every multiple of 16 up to 512,
and pads each step to the next captured size. Steps carry 17, 100, 300 and 600
tokens --- decode tokens, or decode tokens mixed with prefill chunks. At what
size does each run, and what share of its slots is padding? Answer again for
graphs at the powers of two. When do the padded slots cost time?

\textbf{Solution.} With vLLM's sizes, 17 tokens run at 24 (29\% padding), 100
at 104 (3.8\%) and 300 at 304 (1.3\%); 600 exceed the largest graph and run
without one, paying their launches but no padding. With powers of two, 17 run
at 32 (47\%), 100 at 128 (22\%) and 300 at 512 (41\%) (all derived). A padded
slot adds arithmetic --- one more row through every matrix product --- but no
weight reads. A decode step of 17 or 100 sequences is far from compute-bound
(decode at 4,096 tokens of context stays bandwidth-bound at every batch that
fits in an H100, \Chref{napkin-math}), so its padded rows cost almost nothing.
A step of 300 tokens that is mostly a prefill chunk is compute-bound
(\Chref{prefill-decode}), and there every padded slot costs its share: vLLM's
1.3\% costs about 1.3\%, while running 300 tokens at 512 would make the step
about 1.7 times as long for no more output. Dense sizes matter where steps are
compute-bound, and the large token counts of mixed batches are exactly where
vLLM's rule is dense.

\subsection{What graphs buy, and what they leave}
\textbf{Problem.} A Llama-1B forward pass on an H100 at batch one launches
about 100 kernels and streams 2.5~GB of BF16 weights. Launches cost
2.1~$\mu$s each, or 1.3~$\mu$s replayed from a CUDA graph
\cite{spector2025megakernel}. (a) Predict the step time with and without
graphs if the kernels themselves streamed at the full 3.35~TB/s. (b) Hazy
Research measured engines of the time at no more than half the H100's
bandwidth and their megakernel at 78\%. How large is the gap between an engine
with graphs and the megakernel, and how much of it do launches explain?

\textbf{Solution.} (a) Streaming takes $2.5 / 3.35 = 0.75$~ms. Eager launches add
0.21~ms, for 0.96~ms; graphs add 0.13~ms, for 0.88~ms --- 8\% faster (derived).
(b) At half the bandwidth, streaming takes 1.49~ms, and with graphs the step is
about 1.62~ms; at 78\% it takes 0.96~ms, the megakernel's ``under a
millisecond'' (derived). The gap is about 0.66~ms, and the graph's launches,
0.13~ms, are a fifth of it. The other four fifths are bandwidth that the
separate kernels leave unused --- by Hazy Research's account, at their
boundaries: multiprocessors idle while a kernel's last wave drains, and every
kernel stalling while its first weights arrive. Graphs remove the issue cost
of a launch; the larger prize is the gap cost, and only removing the
boundaries, by fusion or a persistent kernel, collects it.

\subsection{When does fusion pay?}
\textbf{Problem.} In the decode step that opened the chapter, llama.cpp fused
each of Llama~3.2~3B's 57 normalizations with the multiplication by its weight,
on activations of 3,072 32-bit floats per token. (a) How many bytes and
dispatches does that save per decoded token, against the 2.0~GB of weights the
step reads? (b) How many bytes does it save in a 512-token prefill, which the
M1 runs at 248 tokens per second (record \texttt{2026-09-23-m1-llama32-3b})? (c)
What kind of fusion would be worth more?

\textbf{Solution.} (a) Unfused, the normalization writes its output and the
multiplication reads it back: $2 \times 3{,}072 \times 4 = 24.6$~KB per
fusion, 1.4~MB per token --- 0.07\% of the weights --- and 57 dispatches of
about 1.3~$\mu$s each, 0.07~ms of a 50~ms step (derived). At decode the
fusion's value is in the dispatches, and even they are small on this machine.
(b) A 512-token batch saves $512 \times 24.6$~KB $= 12.6$~MB per fusion,
0.72~GB in all; at the M1's measured 46--50~GB/s that is about 15~ms of a
2.07~s prefill, 0.7\% (derived). (c) Fusion pays in proportion to the traffic
of the intermediate it removes, measured against the step's total. A
normalization's output is one vector per token, while a projection reads a
whole matrix; the fusion that changed inference is the one whose intermediate
grows fastest --- FlashAttention's, which never writes the score matrix, whose
size grows with the square of the context (\Chref{flashattention}). Fusions of
a matrix product's epilogue, such as dequantizing weights inside the product
(\Chref{weight-quantization}) or quantizing an activation as the previous
kernel writes it, follow the same rule: they remove a full pass over a large
tensor.

